% Options for packages loaded elsewhere
% Options for packages loaded elsewhere
\PassOptionsToPackage{unicode}{hyperref}
\PassOptionsToPackage{hyphens}{url}
%
\documentclass[
  english,
  russian,
  12pt,
  a4paper,
  DIV=11,
  numbers=noendperiod]{scrreprt}
\usepackage{xcolor}
\usepackage{amsmath,amssymb}
\setcounter{secnumdepth}{5}
\usepackage{iftex}
\ifPDFTeX
  \usepackage[T1]{fontenc}
  \usepackage[utf8]{inputenc}
  \usepackage{textcomp} % provide euro and other symbols
\else % if luatex or xetex
  \usepackage{unicode-math} % this also loads fontspec
  \defaultfontfeatures{Scale=MatchLowercase}
  \defaultfontfeatures[\rmfamily]{Ligatures=TeX,Scale=1}
\fi
\usepackage{lmodern}
\ifPDFTeX\else
  % xetex/luatex font selection
\fi
% Use upquote if available, for straight quotes in verbatim environments
\IfFileExists{upquote.sty}{\usepackage{upquote}}{}
\IfFileExists{microtype.sty}{% use microtype if available
  \usepackage[]{microtype}
  \UseMicrotypeSet[protrusion]{basicmath} % disable protrusion for tt fonts
}{}
\usepackage{setspace}
% Make \paragraph and \subparagraph free-standing
\makeatletter
\ifx\paragraph\undefined\else
  \let\oldparagraph\paragraph
  \renewcommand{\paragraph}{
    \@ifstar
      \xxxParagraphStar
      \xxxParagraphNoStar
  }
  \newcommand{\xxxParagraphStar}[1]{\oldparagraph*{#1}\mbox{}}
  \newcommand{\xxxParagraphNoStar}[1]{\oldparagraph{#1}\mbox{}}
\fi
\ifx\subparagraph\undefined\else
  \let\oldsubparagraph\subparagraph
  \renewcommand{\subparagraph}{
    \@ifstar
      \xxxSubParagraphStar
      \xxxSubParagraphNoStar
  }
  \newcommand{\xxxSubParagraphStar}[1]{\oldsubparagraph*{#1}\mbox{}}
  \newcommand{\xxxSubParagraphNoStar}[1]{\oldsubparagraph{#1}\mbox{}}
\fi
\makeatother


\usepackage{longtable,booktabs,array}
\newcounter{none} % for unnumbered tables
\usepackage{calc} % for calculating minipage widths
% Correct order of tables after \paragraph or \subparagraph
\usepackage{etoolbox}
\makeatletter
\patchcmd\longtable{\par}{\if@noskipsec\mbox{}\fi\par}{}{}
\makeatother
% Allow footnotes in longtable head/foot
\IfFileExists{footnotehyper.sty}{\usepackage{footnotehyper}}{\usepackage{footnote}}
\makesavenoteenv{longtable}
\usepackage{graphicx}
\makeatletter
\newsavebox\pandoc@box
\newcommand*\pandocbounded[1]{% scales image to fit in text height/width
  \sbox\pandoc@box{#1}%
  \Gscale@div\@tempa{\textheight}{\dimexpr\ht\pandoc@box+\dp\pandoc@box\relax}%
  \Gscale@div\@tempb{\linewidth}{\wd\pandoc@box}%
  \ifdim\@tempb\p@<\@tempa\p@\let\@tempa\@tempb\fi% select the smaller of both
  \ifdim\@tempa\p@<\p@\scalebox{\@tempa}{\usebox\pandoc@box}%
  \else\usebox{\pandoc@box}%
  \fi%
}
% Set default figure placement to htbp
\def\fps@figure{htbp}
\makeatother



\ifLuaTeX
\usepackage[bidi=basic,shorthands=off,provide=*]{babel}
\else
\usepackage[bidi=default,shorthands=off,provide=*]{babel}
\fi
\ifLuaTeX
  \usepackage{selnolig} % disable illegal ligatures
\fi


\setlength{\emergencystretch}{3em} % prevent overfull lines

\providecommand{\tightlist}{%
  \setlength{\itemsep}{0pt}\setlength{\parskip}{0pt}}



 
\usepackage[style=gost-numeric,backend=biber,langhook=extras,autolang=other*]{biblatex}
\addbibresource{bib/cite.bib}

\usepackage[]{csquotes}

\usepackage{indentfirst}
\usepackage{float}
\floatplacement{figure}{H}
\IfFileExists{plex-otf.sty}{
  %% Full TeXlive
  \usepackage[
    % math,
    RM={Scale=0.94},SS={Scale=0.94},SScon={Scale=0.94},TT={Scale=MatchLowercase,FakeStretch=0.9},
    DefaultFeatures={Ligatures=Common}
  ]{plex-otf}
  % \usepackage[Scale=MatchUppercase]{juliamono}
}{
  %% TinyTeX
  \usepackage{libertine}
}
\KOMAoption{captions}{tableheading}
\makeatletter
\@ifpackageloaded{caption}{}{\usepackage{caption}}
\AtBeginDocument{%
\ifdefined\contentsname
  \renewcommand*\contentsname{Содержание}
\else
  \newcommand\contentsname{Содержание}
\fi
\ifdefined\listfigurename
  \renewcommand*\listfigurename{Список иллюстраций}
\else
  \newcommand\listfigurename{Список иллюстраций}
\fi
\ifdefined\listtablename
  \renewcommand*\listtablename{Список таблиц}
\else
  \newcommand\listtablename{Список таблиц}
\fi
\ifdefined\figurename
  \renewcommand*\figurename{Рисунок}
\else
  \newcommand\figurename{Рисунок}
\fi
\ifdefined\tablename
  \renewcommand*\tablename{Таблица}
\else
  \newcommand\tablename{Таблица}
\fi
}
\@ifpackageloaded{float}{}{\usepackage{float}}
\floatstyle{ruled}
\@ifundefined{c@chapter}{\newfloat{codelisting}{h}{lop}}{\newfloat{codelisting}{h}{lop}[chapter]}
\floatname{codelisting}{Список}
\newcommand*\listoflistings{\listof{codelisting}{Листинги}}
\makeatother
\makeatletter
\makeatother
\makeatletter
\@ifpackageloaded{caption}{}{\usepackage{caption}}
\@ifpackageloaded{subcaption}{}{\usepackage{subcaption}}
\makeatother
\usepackage{bookmark}
\IfFileExists{xurl.sty}{\usepackage{xurl}}{} % add URL line breaks if available
\urlstyle{same}
\hypersetup{
  pdftitle={Лабораторная работа №1},
  pdfauthor={Лаптев Тимофей Сергееви},
  pdflang={ru-RU},
  hidelinks,
  pdfcreator={LaTeX via pandoc}}


\title{Лабораторная работа №1}
\usepackage{etoolbox}
\makeatletter
\providecommand{\subtitle}[1]{% add subtitle to \maketitle
  \apptocmd{\@title}{\par {\large #1 \par}}{}{}
}
\makeatother
\subtitle{Установка и конфигурация операционной системы на виртуальную
машину}
\author{Лаптев Тимофей Сергееви}
\date{}
\begin{document}
\maketitle

\renewcommand*\contentsname{Содержание}
{
\setcounter{tocdepth}{1}
\tableofcontents
}
\listoffigures
\listoftables

\setstretch{1.5}
\chapter{Цель
работы}\label{ux446ux435ux43bux44c-ux440ux430ux431ux43eux442ux44b}

Целью данной работы является приобретение практических навыков установки
операционной системы на виртуальную машину, настройки минимально
необходимых для дальнейшей работы сервисов.

\chapter{Задание}\label{ux437ux430ux434ux430ux43dux438ux435}

\begin{itemize}
\tightlist
\item
  Установить Linux (Rocky) на виртуальную машину VirtualBox.
\item
  Выполнить базовую конфигурацию системы и сетевых параметров.
\item
  Подготовить систему для выполнения последующих лабораторных работ.
\end{itemize}

\chapter{Теоретическое
введение}\label{ux442ux435ux43eux440ux435ux442ux438ux447ux435ux441ux43aux43eux435-ux432ux432ux435ux434ux435ux43dux438ux435}

Виртуализация позволяет запускать гостевые операционные системы в
изолированной среде на одном физическом компьютере. Для учебных задач
удобно использовать VirtualBox, который поддерживает создание
виртуальных дисков, настройку ресурсов и подключение ISO-образов для
установки ОС. После установки важно выполнить базовую настройку системы:
задать имя хоста, создать пользователя, настроить сеть и убедиться в
корректной работе окружения.

\chapter{Выполнение лабораторной
работы}\label{ux432ux44bux43fux43eux43bux43dux435ux43dux438ux435-ux43bux430ux431ux43eux440ux430ux442ux43eux440ux43dux43eux439-ux440ux430ux431ux43eux442ux44b}

\begin{enumerate}
\def\labelenumi{\arabic{enumi}.}
\item
  Перешёл в каталог для временных данных и создал рабочую директорию: cd
  /var/tmp и mkdir /var/tmp/.

  Место для изображения: создание каталога в терминале.
\item
  Запустил VirtualBox и проверил, что папка для виртуальных машин
  указывает на /var/tmp/. Место для изображения: окно настроек
  VirtualBox с путём к папке виртуальных машин.
\item
  Создал виртуальную машину с именем логина, типом Linux RedHat
  (64-bit), объёмом RAM 2048 МБ. Место для изображения: шаги мастера
  создания ВМ с выбором типа ОС и объёма памяти.
\item
  Настроил виртуальный диск VDI (динамический) объёмом 40 ГБ и подключил
  ISO-образ Rocky Linux. Место для изображения: выбор формата диска и
  подключение ISO-образа.
\item
  Выполнил установку системы: выбрал язык интерфейса, часовой пояс,
  раскладки клавиатуры, окружение Server with GUI и дополнение
  Development Tools, отключил KDUMP. Место для изображения: экраны
  установки с выбором языка и параметров установки.
\item
  Включил сетевое соединение, задал имя узла user.localdomain, установил
  пароли для root и пользователя с правами администратора. Место для
  изображения: настройки сети и имя узла, создание пользователя.
\item
  После установки принял лицензионные условия, извлёк ISO-образ и
  перезагрузил систему. Место для изображения: экран принятия лицензии и
  завершения установки.
\item
  Подключил образ дополнений гостевой ОС через меню устройств и
  установил их, затем перезагрузил систему. Место для изображения:
  подключение Guest Additions.
\item
  Проверил соответствие имени пользователя и хоста соглашению об
  именовании: adduser -G wheel , passwd , hostnamectl set-hostname ,
  hostnamectl. Место для изображения: вывод команды hostnamectl.
\end{enumerate}

\chapter{Домашнее
задание}\label{ux434ux43eux43cux430ux448ux43dux435ux435-ux437ux430ux434ux430ux43dux438ux435}

Проанализировал вывод загрузки системы и определил ключевые параметры с
помощью dmesg и grep: версия ядра, частота и модель процессора, объём
доступной памяти, тип гипервизора, тип файловой системы корневого
раздела и порядок монтирования. Место для изображения: фрагменты вывода
dmesg с найденными параметрами.

\chapter{Выводы}\label{ux432ux44bux432ux43eux434ux44b}

В ходе работы установил Rocky Linux в VirtualBox, выполнил первичную
настройку виртуальной машины, задал корректные учётные данные и проверил
параметры загрузки системы.

\chapter*{Список
литературы}\label{ux441ux43fux438ux441ux43eux43a-ux43bux438ux442ux435ux440ux430ux442ux443ux440ux44b}
\addcontentsline{toc}{chapter}{Список литературы}

\printbibliography[heading=none]





\end{document}
